\documentclass[12pt,a4paper]{article}

\usepackage[margin=2.5cm]{geometry}
\usepackage{booktabs}
\usepackage{amsmath}
\usepackage{amssymb}
\usepackage{threeparttable}
\usepackage{setspace}
\usepackage{hyperref}
\usepackage{xcolor}
\usepackage{array}
\usepackage{caption}

\onehalfspacing

\newcolumntype{C}[1]{>{\centering\arraybackslash}p{#1}}

\title{\textbf{Sample Selection in the Estimation Sample}}
\date{\today}

\begin{document}

\maketitle

\section*{Overview}

Our main estimation sample requires workers to have completed \emph{at least two full
spell cycles} (unstable $\to$ stable $\to$ unstable) within the observation window. This
mechanical requirement introduces sample selection: workers who cycle quickly are
over-represented relative to the working population.

This note describes a back-of-envelope exercise designed to \emph{bound} how severe
this selection is. We do not claim to solve the selection problem; rather, we
characterise it quantitatively so we can be precise about the external validity of our
estimates.

\section{The Model}

We assume a \textbf{discrete $2\times 2$ distribution} of individual permanent
heterogeneity in spell exit hazards:
\begin{itemize}
    \item $\lambda_i \in \{\lambda_H, \lambda_L\}$: hazard out of \emph{unstable}
    spells.
    \item $\delta_i \in \{\delta_H, \delta_L\}$: hazard out of \emph{stable}
    spells.
\end{itemize}
This gives four worker types with population shares $\pi_{jk}$:

\begin{center}
\begin{tabular}{lll}
\toprule
Type & $(\lambda,\delta)$ & Interpretation\\
\midrule
HH & $(\lambda_H, \delta_H)$ & Fast cycler (high frequency in both directions)\\
HL & $(\lambda_H, \delta_L)$ & Stable worker (exits unstable quickly, stays stable long)\\
LH & $(\lambda_L, \delta_H)$ & Unstable worker (stuck in unstable, exits stable quickly)\\
LL & $(\lambda_L, \delta_L)$ & Slow cycler (low frequency in both directions)\\
\bottomrule
\end{tabular}
\end{center}

\noindent
Spell durations are exponential, so the career path within any observation window of
length $T$ is a pure-birth continuous-time Markov chain alternating between unstable
and stable states. 
\textcolor{blue}{If spell durations are exponential: Exponential distribution is continuous and memoryless. So, the future evolution of a worker's career depends only on current state, not on how long they have been in that state, or on any history, which is the Markov property, hence why a Markov Chain matches.}
The worker's state at the end of the window determines their
\textbf{bin}:

\begin{center}
\begin{tabular}{cl}
\toprule
Bin & Spells completed at horizon $T$\\
\midrule
0 & Still in first stable spell (entered stable, never exited)\\
0.5 & One started but unfinished unstable spell\\
1 & One complete cycle (U $\to$ S, ongoing)\\
1.5 & One complete cycle plus one unfinished unstable spell\\
2 & Two complete cycles (estimation sample)\\
\bottomrule
\end{tabular}
\end{center}

\noindent
Workers in bin~2 are our estimation sample. Because HH types cycle faster, they
accumulate two complete cycles more rapidly and are therefore mechanically
over-represented in bin~2.

\subsection*{Closed-form bin probabilities}

For each type $(\lambda_j, \delta_k)$ at horizon $T$, the probability of falling in
each bin is available in closed form as a function of survival probabilities of sums of
independent exponential random variables. The derivations are collected in the companion
technical note (\texttt{derivation\_survival\_functions.tex}). \\ \textcolor{blue}{To land in a given bin, a worker must complete a specific sequence of spells before the horizon T. For example, bin 2 (two complete cycles) requires completing four spell exits. The time to complete those four exits is the sum of four independent exponential random variables, two with rate $\lambda$ and two with rate $\delta$ (since spell durations are independent exponentials by assumption).
The probability of being in bin 2 at T is therefore the probability that this sum is less than T, which is the CDF of a sum of independent exponentials evaluated at T (or, one minus the survival function). Sums of independent exponentials with different rates follow a hypoexponential distribution, which has a known closed-form expression. So for each bin, we write down which sequence of spells is needed, sum the corresponding exponentials, and evaluate the closed-form survival function at T.}

\section{Estimation}

We observe the empirical bin shares at four observation windows:

\begin{center}
\begin{tabular}{lrrrrrr}
\toprule
Window & $T$ (yr) & Bin 0 & Bin 0.5 & Bin 1 & Bin 1.5 & Bin 2\\
\midrule
2011--2021 & 10 & 0.338 & 0.214 & 0.256 & 0.069 & 0.123\\
2005--2021 & 16 & 0.266 & 0.171 & 0.277 & 0.076 & 0.211\\
1999--2021 & 22 & 0.219 & 0.151 & 0.291 & 0.075 & 0.264\\
1992--2021 & 29 & 0.206 & 0.135 & 0.266 & 0.069 & 0.326\\
\bottomrule
\end{tabular}
\end{center}

\textcolor{magenta}{\small{We are assuming that the Bin 0 workers have a non-zero probability (\delta) to end up in a temp contract or unemployed.}}
\\
\noindent
We estimate the seven parameters $(\lambda_H, \lambda_L, \delta_H, \delta_L, \pi_{HH},
\pi_{HL}, \pi_{LH})$ (with $\pi_{LL} = 1 - \pi_{HH} - \pi_{HL} - \pi_{LH}$) by
\textbf{GMM} (Generalised Method of Moments), minimising a weighted sum of squared deviations between model-implied and
empirical bin shares across all four windows (20 moments). Weights are the inverse
sampling variances of each share.

The model moments are analytical (no simulation), so the objective is smooth. We use
differential evolution (global search) followed by local gradient-based polishing to
avoid local minima.

\paragraph{Identification remark.} The hazard $\delta_H$ (exit rate from stable spells
for H-$\delta$ types) is weakly identified: at the estimated $\delta_H \approx 0.041$,
the survival probability $e^{-\delta_H T}$ is already negligible for all four horizons,
so the data cannot separately identify further increases. The parameter $\delta_H$
should be interpreted as a lower bound on the true high-delta exit rate. Importantly,
the type \emph{shares} $\pi_{HH} + \pi_{HL}$ (high-$\lambda$ workers) are much more
robustly identified than the individual shares, as they do not hinge on the delta
dimension. \textcolor{blue}{The total share is identified because it depends on $\lambda_H$ which is well identified so the uncertainty about $\delta_H$ only affects the split within the high-$\lambda$ group.}

\section{Parameter Estimates}

\begin{center}
\begin{tabular}{llrr}
\toprule
Parameter & & Estimate & Implied mean dur. (months)\\
\midrule
$\lambda_H$ & Unstable exit (H type) & 0.0258 & 38.8\\
$\lambda_L$ & Unstable exit (L type) & 0.0015 & 667.0\\
$\delta_H$  & Stable exit (H type)   & 0.0409 & 24.4\\
$\delta_L$  & Stable exit (L type)   & 0.0022 & 445.4\\
\midrule
$\pi_{HH}$ & Fast cycler      & 0.156 &\\
$\pi_{HL}$ & Stable worker    & 0.651 &\\
$\pi_{LH}$ & Unstable worker  & 0.180 &\\
$\pi_{LL}$ & Slow cycler      & 0.012 &\\
\bottomrule
\end{tabular}
\end{center}

\noindent
The dominant type is the \textbf{stable worker} (HL): 65\% of workers exit unstable spells quickly but remain in stable jobs for a long time (confirming that \textit{stable} and \textit{unstable} are meaningful names). The high-$\lambda$
group as a whole ($\pi_{HH} + \pi_{HL} = 0.807$) comprises over 80\% of the
population. Positive-correlation types (HH and LL, both fast or both slow) account for
only 17\% of workers; negative-correlation types dominate.

\section{Selection Results}
\label{sec:results}

Table~\ref{tab:composition} reports the type composition of the population, the
one-cycle sample (workers in bins 1 or 1.5), and the two-cycle estimation sample
(bin 2), at each observation window.

\begin{table}[h!]
\centering
\caption{Type composition by sample group and observation window}
\label{tab:composition}
\begin{threeparttable}
\begin{tabular}{llC{1.2cm}C{1.2cm}C{1.2cm}C{1.2cm}C{1.6cm}C{1.6cm}C{1.2cm}}
\toprule
Window & Group & HH & HL & LH & LL
       & $\mathbb{E}[\lambda]$ & $\mathbb{E}[\delta]$ & Share\\
\midrule
\multirow{4}{*}{10 yr}
 & Population        & 0.156 & 0.651 & 0.180 & 0.012 & 0.0211 & 0.0153 & 1.000\\
 & 1 completed spell & 0.141 & 0.787 & 0.069 & 0.003 & 0.0240 & 0.0104 & 0.360\\
 & 2+ completed spells & 0.741 & 0.248 & 0.011 & 0.000 & 0.0255 & 0.0313 & 0.125\\[3pt]
 & \multicolumn{1}{l}{\textit{Ratio (2+ spells / Pop.)}}
   & \textit{4.73$\times$} & \textit{0.38$\times$} & \textit{0.06$\times$} & \textit{0.00$\times$} & & & \\
\midrule
\multirow{4}{*}{16 yr}
 & Population          & 0.156 & 0.651 & 0.180 & 0.012 & 0.0211 & 0.0153 & 1.000\\
 & 1 completed spell   & 0.047 & 0.842 & 0.106 & 0.004 & 0.0231 & 0.0082 & 0.357\\
 & 2+ completed spells & 0.645 & 0.335 & 0.020 & 0.000 & 0.0253 & 0.0280 & 0.213\\[3pt]
 & \multicolumn{1}{l}{\textit{Ratio (2+ spells / Pop.)}}
   & \textit{4.13$\times$} & \textit{0.51$\times$} & \textit{0.11$\times$} & \textit{0.01$\times$} & & & \\
\midrule
\multirow{4}{*}{22 yr}
 & Population          & 0.156 & 0.651 & 0.180 & 0.012 & 0.0211 & 0.0153 & 1.000\\
 & 1 completed spell   & 0.012 & 0.847 & 0.136 & 0.006 & 0.0223 & 0.0080 & 0.353\\
 & 2+ completed spells & 0.551 & 0.418 & 0.030 & 0.000 & 0.0250 & 0.0247 & 0.276\\[3pt]
 & \multicolumn{1}{l}{\textit{Ratio (2+ spells / Pop.)}}
   & \textit{3.52$\times$} & \textit{0.64$\times$} & \textit{0.17$\times$} & \textit{0.02$\times$} & & & \\
\midrule
\multirow{4}{*}{29 yr}
 & Population          & 0.156 & 0.651 & 0.180 & 0.012 & 0.0211 & 0.0153 & 1.000\\
 & 1 completed spell   & 0.002 & 0.828 & 0.162 & 0.008 & 0.0216 & 0.0086 & 0.349\\
 & 2+ completed spells & 0.462 & 0.495 & 0.043 & 0.000 & 0.0247 & 0.0218 & 0.337\\[3pt]
 & \multicolumn{1}{l}{\textit{Ratio (2+ spells / Pop.)}}
   & \textit{2.95$\times$} & \textit{0.76$\times$} & \textit{0.24$\times$} & \textit{0.03$\times$} & & & \\
\bottomrule
\end{tabular}
\begin{tablenotes}[flushleft]\footnotesize
\item \textit{Notes:} Type shares are $P(\text{type} \mid \text{group})$ computed via
Bayes' rule from the estimated type distribution and type-specific bin probabilities.
$\mathbb{E}[\lambda]$ and $\mathbb{E}[\delta]$ are the group-conditional means of the
unstable and stable spell exit hazards. ``Share'' is the fraction of all workers
falling in each group at the given horizon. Hazard rates are per month.
\end{tablenotes}
\end{threeparttable}
\end{table}

\section{Main Findings}

\paragraph{Severe selection toward fast cyclers.}
At the 10-year window, fast cyclers (HH) make up 74\% of the estimation sample despite
representing only 16\% of the population---a 4.7-fold over-representation. Unstable
workers (LH) are almost entirely absent: they fall to 1\% of the estimation sample
(population share: 18\%).

\paragraph{Selection persists at longer windows.}
Even over a 29-year window, HH workers remain nearly 3$\times$ over-represented
(2.95$\times$), and LH workers are still only 4\% of the two-cycle sample. Extending
the observation window substantially improves coverage but does not eliminate the
selection problem within empirically feasible time spans.

\paragraph{The average exit hazard is close to the population one.}
At the 10-year window, the conditional mean unstable-spell exit rate in the estimation
sample is $\bar\lambda = 0.0255$ per month vs.\ $0.0211$ in the population, because the high $\lambda$ types dominate in the economy anyway. The conditional mean stable-spell exit rate is $\bar\delta = 0.0313$ vs.\ $0.0153$
($+105\%$). Selection on $\delta$ is far more severe than selection on $\lambda$.

\paragraph{High-$\lambda$ workers are over-represented, as expected.}
The share of high-$\lambda$ types ($\pi_{HH} + \pi_{HL} = 0.807$ in population) rises
to over 98\% of the two-cycle sample at 10 years and remains above 95\% at 29 years.
This feature is robust to the exact identification of $\delta_H$: whatever the true
split between HH and HL within the high-$\lambda$ group, the group itself is almost
universally over-represented in the estimation sample.

\paragraph{Implications for external validity.}
Our main estimates are identified primarily off high-$\lambda$ workers (those who exit
unstable cycles quickly, one way or another). Workers with persistently low
$\lambda$---who remain stuck in unstable cycles for long stretches---are largely
absent from our sample, but our calculations show that they are not a sizeable part of the population. However, it is good to keep in mind that there is a group that we cannot catch easily the ones with \textit{very} long unstable spells. This puts a limit to how much heterogeneity we can find, but it is unavoidable given the finite lifetime of people.\\

\textit{Note: Cyntia Doniger suggested a Lamadon and Manresa type solution of "clustering" or grouping people with one completed spell only based on observables to increase the sample. That is, getting people with one spell paired up based on characteristics.}

\section{Implications for the Variance Decomposition}

A key motivation for quantifying sample selection is understanding what our variance
decomposition captures. The decomposition partitions the total variance of unstable spell
durations into a \emph{between-worker heterogeneity} component ($\sigma_b^2$) and an
\emph{excess variance} component ($\sigma_e^2$, the paper's measure of duration
dependence). We ask: what would the decomposition return in a world where the data are
generated by the estimated model---that is, a world with \emph{no duration dependence}
by construction?

\paragraph{Simulation design.} Using the type composition of the two-cycle estimation
sample (bin~2) at each horizon, we simulate two sequential unstable spell durations for
each worker using pure exponential distributions. High-$\lambda$ workers (types HH and
HL, joint share $> 96\%$ in all windows) draw from $\text{Exp}(\lambda_H)$; low-$\lambda$
workers from $\text{Exp}(\lambda_L)$. We then apply the paper's variance decomposition to
the simulated sample.

\paragraph{Results.} Table~\ref{tab:vd_sim} reports the shares of total variance
attributed to duration dependence and heterogeneity under this benchmark.

\begin{table}[h!]
\centering
\caption{Variance decomposition in the simulated estimation sample (no duration dependence)}
\label{tab:vd_sim}
\begin{tabular}{lrrrr}
\toprule
Window & $T$ (yr) & High-$\lambda$ share & DD share ($\sigma_e^2/\sigma_y^2$) & HT share ($\sigma_b^2/\sigma_y^2$)\\
\midrule
2011--2021 & 10 & 0.99 & 0.023 & 0.427\\
2005--2021 & 16 & 0.98 & $-$0.006 & 0.453\\
1999--2021 & 22 & 0.97 & 0.009 & 0.453\\
1992--2021 & 29 & 0.96 & $-$0.008 & 0.463\\
\bottomrule
\end{tabular}
\end{table}

\paragraph{Interpretation.} Two findings stand out.

\emph{Duration dependence is essentially zero.} The excess variance share fluctuates
between $-0.008$ and $0.023$ across windows, consistent with sampling noise around
zero. This is the expected result: spell durations are exponential by construction, so
no duration dependence is present. The decomposition recovers this correctly, providing
a validation that the estimator does not spuriously attribute variance to duration
dependence.

\emph{Heterogeneity is non-trivial despite near-homogeneity in $\lambda$.} Even though
over 96--99\% of the estimation sample are high-$\lambda$ workers (HH or HL), the
heterogeneity share is 43--46\%. The mechanism is the minority of low-$\lambda$ workers
(1--4\% of the sample) whose implied mean unstable-spell duration is approximately 667
months---orders of magnitude longer than the 39 months for high-$\lambda$ types. These
few workers are outliers with very long observed spells, which inflate the between-worker
variance substantially. The resulting heterogeneity share would therefore appear even in
the absence of any genuine dispersion in permanent characteristics among the majority of
the sample.

\documentclass[12pt,a4paper]{article}

\usepackage[margin=2.5cm]{geometry}
\usepackage{booktabs}
\usepackage{amsmath}
\usepackage{amssymb}
\usepackage{threeparttable}
\usepackage{setspace}
\usepackage{hyperref}
\usepackage{xcolor}
\usepackage{array}
\usepackage{caption}

\onehalfspacing

\newcolumntype{C}[1]{>{\centering\arraybackslash}p{#1}}

\title{\textbf{Sample Selection in the Estimation Sample}}
\date{\today}

\begin{document}

\maketitle

\section*{Overview}

Our main estimation sample requires workers to have completed \emph{at least two full
spell cycles} (unstable $\to$ stable $\to$ unstable) within the observation window. This
mechanical requirement introduces sample selection: workers who cycle quickly are
over-represented relative to the working population.

This note describes a back-of-envelope exercise designed to \emph{bound} how severe
this selection is. We do not claim to solve the selection problem; rather, we
characterise it quantitatively so we can be precise about the external validity of our
estimates.

\section{The Model}

We assume a \textbf{discrete $2\times 2$ distribution} of individual permanent
heterogeneity in spell exit hazards:
\begin{itemize}
    \item $\lambda_i \in \{\lambda_H, \lambda_L\}$: hazard out of \emph{unstable}
    spells.
    \item $\delta_i \in \{\delta_H, \delta_L\}$: hazard out of \emph{stable}
    spells.
\end{itemize}
This gives four worker types with population shares $\pi_{jk}$:

\begin{center}
\begin{tabular}{lll}
\toprule
Type & $(\lambda,\delta)$ & Interpretation\\
\midrule
HH & $(\lambda_H, \delta_H)$ & Fast cycler (high frequency in both directions)\\
HL & $(\lambda_H, \delta_L)$ & Stable worker (exits unstable quickly, stays stable long)\\
LH & $(\lambda_L, \delta_H)$ & Unstable worker (stuck in unstable, exits stable quickly)\\
LL & $(\lambda_L, \delta_L)$ & Slow cycler (low frequency in both directions)\\
\bottomrule
\end{tabular}
\end{center}

\noindent
Spell durations are exponential, so the career path within any observation window of
length $T$ is a pure-birth continuous-time Markov chain alternating between unstable
and stable states. 
\textcolor{blue}{If spell durations are exponential: Exponential distribution is continuous and memoryless. So, the future evolution of a worker's career depends only on current state, not on how long they have been in that state, or on any history, which is the Markov property, hence why a Markov Chain matches.}
The worker's state at the end of the window determines their
\textbf{bin}:

\begin{center}
\begin{tabular}{cl}
\toprule
Bin & Spells completed at horizon $T$\\
\midrule
0 & Still in first stable spell (entered stable, never exited)\\
0.5 & One started but unfinished unstable spell\\
1 & One complete cycle (U $\to$ S, ongoing)\\
1.5 & One complete cycle plus one unfinished unstable spell\\
2 & Two complete cycles (estimation sample)\\
\bottomrule
\end{tabular}
\end{center}

\noindent
Workers in bin~2 are our estimation sample. Because HH types cycle faster, they
accumulate two complete cycles more rapidly and are therefore mechanically
over-represented in bin~2.

\subsection*{Closed-form bin probabilities}

For each type $(\lambda_j, \delta_k)$ at horizon $T$, the probability of falling in
each bin is available in closed form as a function of survival probabilities of sums of
independent exponential random variables. The derivations are collected in the companion
technical note (\texttt{derivation\_survival\_functions.tex}). \\ \textcolor{blue}{To land in a given bin, a worker must complete a specific sequence of spells before the horizon T. For example, bin 2 (two complete cycles) requires completing four spell exits. The time to complete those four exits is the sum of four independent exponential random variables, two with rate $\lambda$ and two with rate $\delta$ (since spell durations are independent exponentials by assumption).
The probability of being in bin 2 at T is therefore the probability that this sum is less than T, which is the CDF of a sum of independent exponentials evaluated at T (or, one minus the survival function). Sums of independent exponentials with different rates follow a hypoexponential distribution, which has a known closed-form expression. So for each bin, we write down which sequence of spells is needed, sum the corresponding exponentials, and evaluate the closed-form survival function at T.}

\section{Estimation}

We observe the empirical bin shares at four observation windows:

\begin{center}
\begin{tabular}{lrrrrrr}
\toprule
Window & $T$ (yr) & Bin 0 & Bin 0.5 & Bin 1 & Bin 1.5 & Bin 2\\
\midrule
2011--2021 & 10 & 0.338 & 0.214 & 0.256 & 0.069 & 0.123\\
2005--2021 & 16 & 0.266 & 0.171 & 0.277 & 0.076 & 0.211\\
1999--2021 & 22 & 0.219 & 0.151 & 0.291 & 0.075 & 0.264\\
1992--2021 & 29 & 0.206 & 0.135 & 0.266 & 0.069 & 0.326\\
\bottomrule
\end{tabular}
\end{center}

\textcolor{magenta}{\small{We are assuming that the Bin 0 workers have a non-zero probability (\delta) to end up in a temp contract or unemployed.}}
\\
\noindent
We estimate the seven parameters $(\lambda_H, \lambda_L, \delta_H, \delta_L, \pi_{HH},
\pi_{HL}, \pi_{LH})$ (with $\pi_{LL} = 1 - \pi_{HH} - \pi_{HL} - \pi_{LH}$) by
\textbf{GMM} (Generalised Method of Moments), minimising a weighted sum of squared deviations between model-implied and
empirical bin shares across all four windows (20 moments). Weights are the inverse
sampling variances of each share.

The model moments are analytical (no simulation), so the objective is smooth. We use
differential evolution (global search) followed by local gradient-based polishing to
avoid local minima.

\paragraph{Identification remark.} The hazard $\delta_H$ (exit rate from stable spells
for H-$\delta$ types) is weakly identified: at the estimated $\delta_H \approx 0.041$,
the survival probability $e^{-\delta_H T}$ is already negligible for all four horizons,
so the data cannot separately identify further increases. The parameter $\delta_H$
should be interpreted as a lower bound on the true high-delta exit rate. Importantly,
the type \emph{shares} $\pi_{HH} + \pi_{HL}$ (high-$\lambda$ workers) are much more
robustly identified than the individual shares, as they do not hinge on the delta
dimension. \textcolor{blue}{The total share is identified because it depends on $\lambda_H$ which is well identified so the uncertainty about $\delta_H$ only affects the split within the high-$\lambda$ group.}

\section{Parameter Estimates}

\begin{center}
\begin{tabular}{llrr}
\toprule
Parameter & & Estimate & Implied mean dur. (months)\\
\midrule
$\lambda_H$ & Unstable exit (H type) & 0.0258 & 38.8\\
$\lambda_L$ & Unstable exit (L type) & 0.0015 & 667.0\\
$\delta_H$  & Stable exit (H type)   & 0.0409 & 24.4\\
$\delta_L$  & Stable exit (L type)   & 0.0022 & 445.4\\
\midrule
$\pi_{HH}$ & Fast cycler      & 0.156 &\\
$\pi_{HL}$ & Stable worker    & 0.651 &\\
$\pi_{LH}$ & Unstable worker  & 0.180 &\\
$\pi_{LL}$ & Slow cycler      & 0.012 &\\
\bottomrule
\end{tabular}
\end{center}

\noindent
The dominant type is the \textbf{stable worker} (HL): 65\% of workers exit unstable spells quickly but remain in stable jobs for a long time (confirming that \textit{stable} and \textit{unstable} are meaningful names). The high-$\lambda$
group as a whole ($\pi_{HH} + \pi_{HL} = 0.807$) comprises over 80\% of the
population. Positive-correlation types (HH and LL, both fast or both slow) account for
only 17\% of workers; negative-correlation types dominate.

\section{Selection Results}
\label{sec:results}

Table~\ref{tab:composition} reports the type composition of the population, the
one-cycle sample (workers in bins 1 or 1.5), and the two-cycle estimation sample
(bin 2), at each observation window.

\begin{table}[h!]
\centering
\caption{Type composition by sample group and observation window}
\label{tab:composition}
\begin{threeparttable}
\begin{tabular}{llC{1.2cm}C{1.2cm}C{1.2cm}C{1.2cm}C{1.6cm}C{1.6cm}C{1.2cm}}
\toprule
Window & Group & HH & HL & LH & LL
       & $\mathbb{E}[\lambda]$ & $\mathbb{E}[\delta]$ & Share\\
\midrule
\multirow{4}{*}{10 yr}
 & Population        & 0.156 & 0.651 & 0.180 & 0.012 & 0.0211 & 0.0153 & 1.000\\
 & 1 completed spell & 0.141 & 0.787 & 0.069 & 0.003 & 0.0240 & 0.0104 & 0.360\\
 & 2+ completed spells & 0.741 & 0.248 & 0.011 & 0.000 & 0.0255 & 0.0313 & 0.125\\[3pt]
 & \multicolumn{1}{l}{\textit{Ratio (2+ spells / Pop.)}}
   & \textit{4.73$\times$} & \textit{0.38$\times$} & \textit{0.06$\times$} & \textit{0.00$\times$} & & & \\
\midrule
\multirow{4}{*}{16 yr}
 & Population          & 0.156 & 0.651 & 0.180 & 0.012 & 0.0211 & 0.0153 & 1.000\\
 & 1 completed spell   & 0.047 & 0.842 & 0.106 & 0.004 & 0.0231 & 0.0082 & 0.357\\
 & 2+ completed spells & 0.645 & 0.335 & 0.020 & 0.000 & 0.0253 & 0.0280 & 0.213\\[3pt]
 & \multicolumn{1}{l}{\textit{Ratio (2+ spells / Pop.)}}
   & \textit{4.13$\times$} & \textit{0.51$\times$} & \textit{0.11$\times$} & \textit{0.01$\times$} & & & \\
\midrule
\multirow{4}{*}{22 yr}
 & Population          & 0.156 & 0.651 & 0.180 & 0.012 & 0.0211 & 0.0153 & 1.000\\
 & 1 completed spell   & 0.012 & 0.847 & 0.136 & 0.006 & 0.0223 & 0.0080 & 0.353\\
 & 2+ completed spells & 0.551 & 0.418 & 0.030 & 0.000 & 0.0250 & 0.0247 & 0.276\\[3pt]
 & \multicolumn{1}{l}{\textit{Ratio (2+ spells / Pop.)}}
   & \textit{3.52$\times$} & \textit{0.64$\times$} & \textit{0.17$\times$} & \textit{0.02$\times$} & & & \\
\midrule
\multirow{4}{*}{29 yr}
 & Population          & 0.156 & 0.651 & 0.180 & 0.012 & 0.0211 & 0.0153 & 1.000\\
 & 1 completed spell   & 0.002 & 0.828 & 0.162 & 0.008 & 0.0216 & 0.0086 & 0.349\\
 & 2+ completed spells & 0.462 & 0.495 & 0.043 & 0.000 & 0.0247 & 0.0218 & 0.337\\[3pt]
 & \multicolumn{1}{l}{\textit{Ratio (2+ spells / Pop.)}}
   & \textit{2.95$\times$} & \textit{0.76$\times$} & \textit{0.24$\times$} & \textit{0.03$\times$} & & & \\
\bottomrule
\end{tabular}
\begin{tablenotes}[flushleft]\footnotesize
\item \textit{Notes:} Type shares are $P(\text{type} \mid \text{group})$ computed via
Bayes' rule from the estimated type distribution and type-specific bin probabilities.
$\mathbb{E}[\lambda]$ and $\mathbb{E}[\delta]$ are the group-conditional means of the
unstable and stable spell exit hazards. ``Share'' is the fraction of all workers
falling in each group at the given horizon. Hazard rates are per month.
\end{tablenotes}
\end{threeparttable}
\end{table}

\section{Main Findings}

\paragraph{Severe selection toward fast cyclers.}
At the 10-year window, fast cyclers (HH) make up 74\% of the estimation sample despite
representing only 16\% of the population---a 4.7-fold over-representation. Unstable
workers (LH) are almost entirely absent: they fall to 1\% of the estimation sample
(population share: 18\%).

\paragraph{Selection persists at longer windows.}
Even over a 29-year window, HH workers remain nearly 3$\times$ over-represented
(2.95$\times$), and LH workers are still only 4\% of the two-cycle sample. Extending
the observation window substantially improves coverage but does not eliminate the
selection problem within empirically feasible time spans.

\paragraph{The average exit hazard is close to the population one.}
At the 10-year window, the conditional mean unstable-spell exit rate in the estimation
sample is $\bar\lambda = 0.0255$ per month vs.\ $0.0211$ in the population, because the high $\lambda$ types dominate in the economy anyway. The conditional mean stable-spell exit rate is $\bar\delta = 0.0313$ vs.\ $0.0153$
($+105\%$). Selection on $\delta$ is far more severe than selection on $\lambda$.

\paragraph{High-$\lambda$ workers are over-represented, as expected.}
The share of high-$\lambda$ types ($\pi_{HH} + \pi_{HL} = 0.807$ in population) rises
to over 98\% of the two-cycle sample at 10 years and remains above 95\% at 29 years.
This feature is robust to the exact identification of $\delta_H$: whatever the true
split between HH and HL within the high-$\lambda$ group, the group itself is almost
universally over-represented in the estimation sample.

\paragraph{Implications for external validity.}
Our main estimates are identified primarily off high-$\lambda$ workers (those who exit
unstable cycles quickly, one way or another). Workers with persistently low
$\lambda$---who remain stuck in unstable cycles for long stretches---are largely
absent from our sample, but our calculations show that they are not a sizeable part of the population. However, it is good to keep in mind that there is a group that we cannot catch easily the ones with \textit{very} long unstable spells. This puts a limit to how much heterogeneity we can find, but it is unavoidable given the finite lifetime of people.\\

\textit{Note: Cyntia Doniger suggested a Lamadon and Manresa type solution of "clustering" or grouping people with one completed spell only based on observables to increase the sample. That is, getting people with one spell paired up based on characteristics.}

\section*{Limitations}

This exercise abstracts from several complications: (i) the discrete $2\times 2$ grid may not be a good approximation of the underlying distribution (there might be more types, for example); (ii) entry into the
labour market is treated as uniform over the observation window; (iii) $\delta_H$ is
only weakly identified (see Section~2), so the exact split between HH and HL types
within the estimation sample is uncertain, though the total high-$\lambda$ share is
stable. The exercise should be read as a plausibility bound, not a precise structural
estimate - that is a paper in itself!

\section{Applying the Variance Decomposition to our estimated population sample}

% A key motivation for quantifying sample selection is understanding what our variance decomposition captures. 
The decomposition partitions the total variance of unstable spell
durations into a \emph{between-worker heterogeneity} component ($\sigma_b^2$) and an
\emph{excess variance} component ($\sigma_e^2$, the paper's measure of duration
dependence). We ask: what would the decomposition return in a world where the data are
generated by the estimated model---that is, a world with \emph{no duration dependence}
by construction?

\paragraph{Simulation design.} Using the type composition of the two-cycle estimation
sample (bin~2) at each horizon, we simulate two sequential unstable spell durations for
each worker. High-$\lambda$ workers (types HH and
HL, joint share $> 96\%$ in all windows) draw from $\text{Exp}(\lambda_H)$; low-$\lambda$
workers from $\text{Exp}(\lambda_L)$, as the sample selection exercise assumes. We then apply the paper's variance decomposition to the simulated sample.

\paragraph{Results.} Table~\ref{tab:vd_sim} reports the shares of total variance attributed to duration dependence and heterogeneity under this benchmark.

\begin{table}[h!]
\centering
\caption{Variance decomposition in the simulated estimation sample}
\label{tab:vd_sim}
\begin{tabular}{lrrrr}
\toprule
Window & $T$ (yr) & High-$\lambda$ share & DD share ($\sigma_e^2/\sigma_y^2$) & HT share ($\sigma_b^2/\sigma_y^2$)\\
\midrule
2011--2021 & 10 & 0.99 & 0.023 & 0.427\\
2005--2021 & 16 & 0.98 & $-$0.006 & 0.453\\
1999--2021 & 22 & 0.97 & 0.009 & 0.453\\
1992--2021 & 29 & 0.96 & $-$0.008 & 0.463\\
\bottomrule
\end{tabular}
\end{table}

\paragraph{Interpretation.} Two findings stand out.

\emph{Duration dependence is essentially zero.} This is by design here, but it echoes our findings in Spain. The excess variance share fluctuates
between $-0.008$ and $0.023$ across windows, consistent with sampling noise around
zero.

\emph{Heterogeneity is very large} Even though
over 96--99\% of the estimation sample are high-$\lambda$ workers (HH or HL), the
heterogeneity share is 43--46\%. This is way above our estimates for Spain. The mechanism here is the minority of low-$\lambda$ workers
(1--4\% of the sample) whose implied mean unstable-spell duration is approximately 667
months---orders of magnitude longer than the 39 months for high-$\lambda$ types. These
few workers are outliers with very long observed spells, which inflate the between-worker variance substantially. The resulting heterogeneity share would therefore appear even in the absence of any genuine dispersion in permanent characteristics among the majority of the sample.

\paragraph{Caveat.} The simulation should not be read as a quantitative prediction of
what the decomposition returns in the actual data. On the duration dependence side the
simulation is consistent with the data: the actual DD shares range from $-$3\% to
$+$5\% across windows (Spain time-window results, main paper), confirming that the
decomposition does not attribute variance to duration dependence. On the heterogeneity
side, however, the simulation substantially \emph{overestimates} the data: the actual
HT shares are 9--16\%, roughly one-third of the simulated 43--46\%. Two limitations of
the 2$\times$2 model drive this gap. First, it collapses all high-$\lambda$ workers
(HH and HL) to a single rate $\lambda_H$, ignoring the within-group dispersion that
is present in the data and that the decomposition captures as heterogeneity. Second, the identification of $\lambda_L$ in a two-point grid is fragile; the estimated value
($\approx667$ months) is extreme and creates variance outliers that mechanically inflate
the heterogeneity share. A richer model with more types would be needed to match the
data's heterogeneity level quantitatively, but that is beyond the scope of this exercise. \textcolor{cyan}{Note to authors: the previous version with continuous types would have avoided this, but then we would be constrained in the $\delta$-$\lambda$ dependency. This is the trade-off we do in the sample selection part.}
The simulation is best interpreted as the worst-case scenario in terms of how bad heterogeneity can get, given the data.



\end{document}


\end{document}
